\documentclass[11pt]{article}

\usepackage[margin=1in]{geometry}
\usepackage{amsmath}
\usepackage{amssymb}
\usepackage{amsthm}
\usepackage[T1]{fontenc}
\usepackage{lmodern}
\usepackage{hyperref}
\hypersetup{colorlinks=true, linkcolor=blue, urlcolor=blue}

% ---- commit hash of the work-in-progress release (STEP 1) ----
\newcommand{\committag}{edd355443952219a6d9f03275bb218feecc2cc4b}

\newcommand{\Aut}{\operatorname{Aut}}
\newcommand{\im}{\operatorname{im}}
\newcommand{\Z}{\mathbb{Z}}
\newcommand{\nubr}{[\nu]}
\newcommand{\Om}{[\Omega]}

\title{$[\nu]\perp[\Omega]$: the $|G|=16$ corner and the moving-$\nu$ cocycle mechanism}
\author{MacBeth}
\date{2026-09-26}

\begin{document}
\maketitle

\begin{center}
\begin{tabular}{rl}
\textbf{For:} & Rick \\
\textbf{Author:} & MacBeth \\
\textbf{Date:} & 2026-09-26 (WAKE deep-work)\\
\textbf{Grade:} & \textsc{Computed} (explicit witness + machine-checked derivation on a witness)\\
\textbf{Registry:} & \texttt{quiver-skew-brace-zs.json} --- node \texttt{cross-independence-nu-omega} (was \texttt{computed})\\
\textbf{WIP commit:} & \texttt{\committag}\\
\end{tabular}
\end{center}

\begin{abstract}
\noindent\textbf{Grade of this note: \textsc{Computed}} (explicit witness + machine-checked
derivation on a witness). \emph{Not proved}: the general converse-as-theorem has one un-closed
step (see \S4).\\[4pt]
Related: \texttt{[[nu-variation-direct-factor-orthogonal]]} (proved 09-25, the invariant-level
orthogonality); \texttt{[[bilinear-skew-brace-h2-identification-proved]]} (peer-reviewed, the
$\Om$ home); \texttt{[[nontrivial-kernel-orthogonal-to-obstruction]]}.
\end{abstract}

\section{Setup}
Skew-brace extension $0\to D\to G\to H\to 0$, $D$ an ideal with $(D,+),(D,\circ)$ abelian.
$a\circ b = a + \lambda_a(b)$, $\lambda:(G,\circ)\to\Aut(G,+)$ a homomorphism.
Section-independent invariants $(\mu,\sigma,\nubr,\Om)$. $L := \im(\lambda|_D)\le\Aut(D,+)$,
residue $[\nu_h]=\nu_h\cdot L\in\Aut(D)/L$.

Already \textsc{proved} (09-25): the $\nu$-orbit is a free transitive $L^{H\setminus 0}$-torsor
(L1); $[\nu_h]$ section-independent (L2); $[\beta],[\bar\tau]$ and the splitting predicate
$\nu$-free \& constant on the torsor (L3); orthogonality holds at the \textbf{invariant} level
(naive gauge-\textsc{orbit} product is \textsc{false} --- flipping $\nu$ forces a coboundary).
Open piece $=$ cross-independence $\nubr\perp\Om$, only \texttt{computed} on $\Z_2\times\Z_4$
(all 6 combos), and only in the $|D|=4$ regime where $L\neq 1\Rightarrow L=\Aut(D)$ is a size
artifact.

\section{The flagged corner --- CLOSED at $|G|=16$ (computed)}
Script: \texttt{scratch/2026-09-26-nu-corner-construction.py} (runs clean).

\medskip\noindent\textbf{Construction.}
$(G,+)=\Z/2\times\Z/8$, $D=\{(0,x)\}\cong\Z/8$, $H=\Z/2$, $\Aut(D)=(\Z/8)^{*}\cong\Z/2\times\Z/2$.
$\lambda$ from the unit-scaling family $\varphi_{u,t}(j,y)=(j,\;u\cdot y + j\cdot t)$,
$u\in\{1,3,5,7\}$, $t\in\{0,4\}$; $U(i,x)=5^{(x\bmod 2)}\cdot 3^{i}\pmod 8$ a homomorphism
$(G,\circ)\to\text{units}$. Two witnesses differing \textbf{only} in the $t$-part (invisible on
$D\Rightarrow$ identical $L$ and $\nu$):
\begin{itemize}
  \item \textbf{G0 (split):} $\lambda^0_{(i,x)}=\varphi_{U(i,x),0}$;
  \item \textbf{G1 (non-split):} $\lambda^1_{(i,x)}=\varphi_{U(i,x),4i}$
        ($T(i,x)=4i$ a homomorphism $(G,\circ)\to\{0,4\}$).
\end{itemize}
Then $L=\{1,5\}$ (order $2$, \textbf{proper nontrivial} in $\Aut(D)$); base lift $s(1)=(1,0)$
gives $\nu_1=\,\cdot\,U(1,0)=\,\cdot\,3\notin L$; residue $[\nu_1]=\{3,7\}\neq L$.

\medskip\noindent\textbf{All six verification items PASS (identically on G0, G1):}
valid skew brace (compat over all $16^3$ triples); $D$ ideal, both ops abelian;
$1\subsetneq L\subsetneq\Aut(D)$; residue nontrivial; L1 ($\nu$-values over the $8$ lifts of
$h=1$ are exactly $\{3,7\}=3\cdot L$, free $L$-torsor), L2, L3;
\textbf{cross-tab:} SAME $(L,\nubr)=(\{1,5\},\{3,7\})$ on both, yet $\Om=0$ for G0
(complement $\{(0,0),(1,0)\}$ closed under both ops) and $\Om\neq 0$ for G1
(both order-$2$ candidates fail $\circ$-closure: $(1,x)\circ(1,x)=(0,4)$).

\medskip\noindent$\Rightarrow$ A witness with \textbf{genuinely moving $\nu$} ($L\neq 1$)
carrying a \textbf{nontrivial residue} ($L\subsetneq\Aut(D)$) exists at $|G|=16$, and
$\nubr\perp\Om$ cross-independence persists there --- no longer a $|D|=4$ artifact.

\section{The moving-$\nu$ cocycle equations (derived natively; machine-checked on a witness)}
Script/derivation: \texttt{scratch/2026-09-26-moving-nu-cocycle-derivation.\{md,py\}} (sympy green).
RY (2601.12371) only cover trivial kernels; ref [20] is RY's own 2102.12235; there is NO
``Letourmy--Vendramin cohomology of skew braces''. So this is native. Cochains valued in $D$;
$\nu_h := \lambda_{s(h)}|_D\in\Aut(D,+)$; $T$ $=$ additive-difference circle factor set.
\begin{itemize}
  \item[\textbf{(a)}] additive $2$-cocycle:
    $\mu_{h_1}\beta(h_2,h_3)+\beta(h_1,h_2+h_3)=\beta(h_1,h_2)+\beta(h_1+h_2,h_3)$, $\mu$ a hom.
  \item[\textbf{(b)}] circle $2$-cocycle ($\circ$-mirror in abelian $(D,\circ)$).
  \item[\textbf{(c)}] cross eq (generalized RY 3.10):
  \begin{align*}
    \mu_h\,\nu_h\,\beta(k,l) + T(h,k+l)
      &= T(h,k) - \mu_{h\circ k}\mu_h^{-1}\beta(h,-h) + \beta(h\circ k,-h)\\
      &\quad{}+ \mu_{(h\circ k)-h}T(h,l) + \beta\big((h\circ k)-h,\; h\circ l\big).
  \end{align*}
  \item action-block \textbf{coupling-1} (no $\beta,\tau$):
    $\nu_h\,\mu_k\,\nu_h^{-1}=\mu_{\lambda^H_h(k)}$, and $\Lambda_{(d,h)}\equiv\nu_h\pmod L$.
\end{itemize}
$\nu$ enters (c) \textbf{exactly once}, as $\mu_h\,\nu_h\,\beta(k,l)$.

\medskip\noindent Machine-checked on the witness $(G,+)=\Z_2\times\Z_4$, $D=\{0\}\times\Z/4$,
$\lambda_{(x,y)}(x',y')=(x',(-1)^{x+y}y')$ [$(D,+)=\Z/4$, $(D,\circ)=$ Klein, $\lambda|_D=$ neg]:
coordinate $\oplus,\circledast$ reproduce true $+,\circ$; (a),(b),(c), coupling-1,
$\Lambda\equiv\nu$ all True over all $H$-triples / $64$ pairs / $3$ sections.

\section{Verdict --- two levels (this reproduces \& explains the 09-25 correction)}
\begin{itemize}
  \item \textbf{Cochain level: GENUINE COUPLING (not $\nu$-free).} The term $\mu_h\,\nu_h\,\beta(k,l)$
    means moving $\nu_h$ within its $L$-coset forces a compensating change in the circle factor
    set $\tau$.
  \item \textbf{Invariant level: INDEPENDENCE.} Because $a\circ b=a+\lambda_a(b)$, once $(G,+)$
    (i.e.\ $\mu,\beta$) and $\lambda$ (i.e.\ $\nu$ and the $L$-part) are fixed, $\circ$ --- hence
    $\tau$ --- is \textbf{determined}, not freely chosen; eq (c) merely \emph{defines} $T$ from
    $\beta$ and the action block, uniquely up to circle-coboundary, and the induced $\tau$
    satisfies (b). So $\nu$ imposes \textbf{no existence obstruction on $\beta$}: for every
    $\beta$ and every admissible $\nu$ a compatible $\tau$ exists $\Rightarrow$ $H^2_{\nubr}$
    (obstruction classes at fixed residue) is $\nubr$-independent as an abstract group; $\nubr$
    is an orthogonal direct factor. The coupling is confined to coboundaries $\to$ invisible to
    $H^{*}$. This is exactly the registered 09-25 shape.
\end{itemize}

\section{The ONE un-closed step (why this is \texttt{computed}, not \texttt{proved})}
The invariant-level independence argument in \S3 rests on: \textbf{eq (c) is auto-solvable for
$T$ given ANY $\beta$, and the resulting $\circ$ is associative / (b) holds automatically.}
Structurally this should follow from ``$\circ$ is determined by $+$ and $\lambda$, and $\lambda$
a homomorphism $\Rightarrow$ $(G,\circ)$ a group'', i.e.\ the RHS of (c) is automatically a
$T$-coboundary. The derivation \textsc{verified} this on the witness and argued it structurally,
but did \textbf{not} close the general ``RHS is always a $T$-coboundary'' step. Until that is a
clean proof (associativity of $\circ$ $\Rightarrow$ solvability of (c) for all $\beta$, uniform
in the admissible residue), the general converse ($\Om\neq 0$ with prescribed $[\nu']$) stays
\textbf{\texttt{computed}} (witnessed at $|D|=4$ and now $|D|=8$), not \texttt{proved}.

\medskip\noindent\textbf{Note on the $|D|=4$ witness.} There $\nu$ and $\beta$
\emph{anti-correlate} ($\nu_1=$ neg $\Leftrightarrow$ $\beta$ in the neg-fixed $2$-torsion), so
the coupling term is present but never ``bites''; separating distinct $\nubr$ needs
$L\subsetneq\Aut(D)$, i.e.\ the $|G|=16$ witness of \S1 --- which is precisely why the corner
mattered.

\section{Next (seed PROVE.md)}
Close \S4: prove that associativity of $a\circ b=a+\lambda_a(b)$ (equivalently $\lambda$ a
skew-brace $\lambda$-map) forces eq (c) to be solvable for $T$ for every additive cocycle
$\beta$, uniformly over admissible residues $\nubr$ --- upgrading
\texttt{cross-independence-nu-omega} from \texttt{computed} to \texttt{proved}. Then the map
$\{\text{extensions/fixed }(\mu,\sigma)\}\to\{\nubr\}\times\{\Om\text{-value}\}$ is a genuine
product (surjective), completing the internal-replacement / decomposition-obstruction
coordinate picture.

\end{document}
